% Ad Atticum 10.1 --- headnote
% ad-atticum-10-01, 3 April 49 BC, Laterium

\textit{Cicero to Atticus, written from Laterium, the
country place of his brother Quintus, on the third day
before the Nones of April 49 BC (the manuscript dateline:
\textit{Scr.\ in Laterio Quinti fratris iii Non.\ Apr.\ a.\
705 (49)}). Cicero has paused on his way from Formiae:
Atticus's letter has reached him, and he says he has
``breathed a little'' for the first time since the
catastrophe began. The pretext for the relief is Atticus's
report that Sextus Peducaeus the younger approves Cicero's
firmness --- which lets Cicero feel as if confirmed by the
elder Sextus, whom he is moved to quote in Greek hexameter
from Hector's last words in the \textit{Iliad}: ``not
without a struggle and inglorious, but having done some
great thing for men of after-time to learn of.'' The body of
the letter is taken up with the same paralysis as the
preceding weeks --- what to do if Caesar's circle invite him
to a peace-council, what to do if he is dispatched as envoy,
whether to be absent ``from this side and from that.''
Solon's law penalising the non-aligned in civil strife is
quoted, and at once set aside.}

\textit{Section 3 carries a textual crux, marked here with
daggers as in the standard text. The Greek tags throng the
letter: \textit{[Greek: t\^on politik\=ot\'at\=on
skemm\'at\=on]} (``of the most properly political
questions'') flags the dilemma about attending a tyrant's
council as a topos of philosophical-political
deliberation; \textit{[Greek: ekph\=on\=esis hup\'ereu]}
records Atticus's habit of breaking in with a ``splendid!''
of approval; \textit{[Greek: \'al\=e]} (``wandering,
drifting'') names the present condition as a kind of
living death; \textit{[Greek: politeut\'eon]} (``one
must play the citizen'') sets up the antithesis that
governs section 4 --- either freely among the wicked or,
with danger, with the honest. The mention of Flavius
receiving a legion and Sicily, of Trebatius as ``an honest
man and citizen,'' and of an unnamed peacemaker who has
sent his son to Brundisium (almost certainly L.\ Caesar)
flesh out a letter whose surface is anxious news from a
brother's villa and whose substance is the recurring
question of what political action is even possible.}
